\documentclass{msukz}	% Подключается файл настроек msukz.cls
% Пожалуйста, не подключайте других пакетов
% и не производите какую-либо настройку.
% Они будут проигнорированы при компиляции сборника.
% При необходимости свяжитесь с техническим секретарем.

\begin{document}
\info % информация о тезисе:
{Касенова А.А.}										% Фамилия И.О. авторов
{Пример оформления тезиса с формулами, рисунками и таблицами}	% Название тезиса
{Бакалавр, студент}											% Статусы (первый с заглавной буквы)
{Казахстанский филиал МГУ имени М.В.Ломоносова}					% Организация
{г.~Астана, Казахстан}											% Город, страна
{assemk2004@gmail.com}									% e-mail (если есть)

В современных задачах кредитного скоринга необходимо выявлять сложные нелинейные зависимости в кредитных историях заемщиков. Классическая логистическая регрессия не всегда способна выявлять такие закономерности. В данной работе проводится исследование использования нейронной сети для прогнозирования дефолта заемщика на основе исторических данных платежной дисциплины. Данные взяты из платформы ODS.ai.

В качестве исходных данных используется матрица X\_pay, состоящая из 25 признаков истории платежей (enc\_paym0,\dots,enc\_paym24). Целевой переменной является флаг дефолта (1 -- дефолт, 0 -- его отсутствие). Для обучения и валидации исходный массив данных, состоящий из 3\,000\,000 клиентов, был разделен на тренировочную (75\%) и тестовую (25\%) выборки.

В работе используется нейронная сеть со следующей структурой:

\onepicture{neural_network}{Структура нейронной сети}

Обученная модель продемонстрировала высокую предсказательную способность. Основные метрики качества модели представлены ниже:

%\begin{center}
%	\begin{longtable}{|c|c|c|c|c|c|c|}
%		\caption{Метрики точности алгоритма Neural Network}
%		\label{metrics_nn} \\
%		\hline
%		\hline
%		\rowcolor[gray]{.9}
%		\specialcell{Flag\\метка} &
%		\specialcell{Precision\\точность} &
%		\specialcell{Recall\\полнота} &
%		\specialcell{f1--score\\f1--мера} &
%		\specialcell{Accuracy\\точность}  &
%		ROC--AUC &
%		Количество    \\
%		\hline
%		\multicolumn{7}{|c|}{Тренировочная выборка}\\
%		\hline
%		0 & 1.00 & 0.97 & 0.99 &
%		\multirow{2}{*}{0.98} &   
%		\multirow{2}{*}{0.98} &  
%		469239 \\                    
%		\cline{1-4} \cline{7-7}
%		1 & 0.95 & 1.00 & 0.97 &  
%		&      & 280761 \\       
%		\hline
%		\multicolumn{7}{|c|}{Тестовая выборка} \\
%		\hline
%		0 & 1.00 & 0.97 & 0.98 &
%		\multirow{2}{*}{0.98} &
%		\multirow{2}{*}{0.98} &
%		156325 \\
%		\cline{1-4} \cline{7-7}
%		1 & 0.95 & 1.00 & 0.97 & & & 93676 \\
%		\hline
%	\end{longtable}
%\end{center}

Метрики $accuracy$ и $ROC-AUC$ демонстрируют одинаковую высокую точность, что также относит нейронные сети к наиболее предпочтительным алгоритмам для прогнозирования дефолта заемщиков.

\onepicture{ROC_AUC_nn_train_test}{Структура нейронной сети}

Обратите внимание на список литературы, который оформлен с помощью команд файла настроек {\tt msukz}. Для лучшего понимания, откройте текст {\tt IvanovAA.tex}.

Не забывайте, что итоговый объем текста не должен превышать 2 страниц (в случае наличия изображений --- 3 страницы). Текст можно сохранять в кодировке {\tt cp1251} (Windows) или {\tt utf8} (Linux).


% список литератутры следует обернуть окружением reference, который задан в файле msukz.cls

% КНИГА: \bookRu, \bookEn, \bookKz
% #1 - авторы
% #2 - название
% #3 - город
% #4 - издательство
% #5 - год
% #6 - страниц

% СТАТЬЯ: \articleRu, \articleEn, \articleKz
% #1 - авторы
% #2 - название
% #3 - журнал
% #4 - год
% #5 - том (необязательно)
% #6 - номер (необязательно)
% #7 - страницы

% ИНТЕРНЕТ: \internet
% #1 - авторы
% #2 - название
% #3 - web-ссылка

% Если ни один из 3 образцов заполнения Вам не подходит, оформите источник самостоятельно:
% ПРОИЗВОЛЬНЫЙ ТЕКСТ: \refer
% #1 - произвольный текст

\begin{references}
\articleRu{Отелбаев М.}{Существование сильного решения уравнения Навье-Стокса}{Математический журнал}{2013}{13}{4}{5--104}
\bookRu{Рябцева Н.Е.}{Язык и естественный интеллект}{М.}{Академия}{2005}{639}
\internet{Gnuplot}{Command-line driven graphing utility}{http://www.gnuplot.info}
\end{references}

\end{document}
